\chapter{Fourier Duality, Wave Packets, and Photon Polarization}

\section{A Periodic Line as a Regulator}

Begin with a line segment of length \(L\), but identify its two ends. A
particle that leaves at one end reappears at the other. This is equivalent
mathematically to motion on a circle, although we need not imagine the
particle literally travelling around a physical ring. The construction is a
regulator: it makes momentum discrete and normalization ordinary, and the
infinite line is recovered by taking \(L\) to infinity.

The endpoint identification requires
\begin{equation}
  \psi(x+L)=\psi(x).
  \label{eq:periodic-line}
\end{equation}
It is convenient to write \(L=2\pi r\). The symbol \(r\) then resembles a
circle radius, but only the circumference \(2\pi r\) matters.

\begin{figure}[tbp]
  \centering
  \includegraphics[width=0.68\textwidth]{lecture_05_figure_02.png}
  \caption{The periodicity condition and the finite interval of length
  \(L\).}
  \lectureframe{00:02:17}
\end{figure}

Momentum eigenfunctions are plane waves. Restoring units before simplifying
the notation,
\begin{equation}
  e^{ipx/\hbar}=e^{ikx},
  \qquad
  k:=\frac{p}{\hbar}.
  \label{eq:k-definition}
\end{equation}
Thus \(k\) is the wave number, while the physical momentum is \(p=\hbar k\).
On one period, a normalized plane wave is
\begin{equation}
  \psi_k(x)=\braket{x}{k}_{\mathrm{box}}
    =\frac{e^{ikx}}{\sqrt{2\pi r}},
  \qquad 0\leq x<2\pi r.
  \label{eq:box-plane-wave}
\end{equation}

\begin{figure}[tbp]
  \centering
  \includegraphics[width=0.58\textwidth]{lecture_05_figure_03.png}
  \caption{The normalized plane wave written in \(p/\hbar\) and \(k\)
  notation.}
  \lectureframe{00:04:20}
\end{figure}

\begin{classroomqa}
\audiencequestion{Why does the normalization contain a square root?
\lecturetimestamp{00:04:29--00:05:26}}
\lecturerresponse{Normalization is imposed on the absolute square. Since
\(\lvert e^{ikx}\rvert^2=1\),
\[
  \int_0^{2\pi r}\lvert e^{ikx}\rvert^2\,dx=2\pi r.
\]
Dividing the amplitude by \(\sqrt{2\pi r}\) divides its absolute square by
\(2\pi r\), leaving unit norm.}
\end{classroomqa}

Periodicity now quantizes the wave number:
\begin{align}
  e^{ik(x+2\pi r)}&=e^{ikx},\\
  e^{i2\pi rk}&=1,\\
  k&=\frac{n}{r},
  \qquad n\in\mathbb Z.
  \label{eq:k-quantization}
\end{align}
The allowed values are separated by
\begin{equation}
  \Delta k=\frac{1}{r}.
  \label{eq:k-spacing}
\end{equation}
As \(r\) grows, the spectrum becomes dense. A dense quantum spectrum is not
by itself a classical limit; it is only the continuum limit of this
kinematical basis.

\section{Orthogonality Before the Continuum}

The momentum operator is Hermitian on periodic functions with the
appropriate domain, so eigenfunctions with different eigenvalues are
orthogonal. Here the theorem can be checked directly. For
\(k_n=n/r\) and \(k_m=m/r\),
\begin{align}
  \braket{k_m}{k_n}_{\mathrm{box}}
  &=\frac{1}{2\pi r}\int_0^{2\pi r}
      e^{i(k_n-k_m)x}\,dx \\
  &=\delta_{mn}.
  \label{eq:box-orthogonality}
\end{align}
When \(m=n\), the integrand is one. When \(m\neq n\), it completes an
integer number of oscillations, and its positive and negative contributions
cancel.

\begin{classroomqa}
\audiencequestion{Is orthogonality here a definition or a result?
\lecturetimestamp{00:11:24--00:12:54}}
\lecturerresponse{It is a result. Distinct eigenvalues of a Hermitian
operator have orthogonal eigenvectors. The integral above verifies that
general theorem for the periodic plane waves; it does not define
orthogonality.}
\end{classroomqa}

The states \(\ket{k_n}_{\mathrm{box}}\) therefore form an orthonormal basis.
At finite \(r\), the correct language is discrete: sums over \(n\) and
Kronecker deltas \(\delta_{mn}\). We now stretch this basis into the
continuum without losing its normalization bookkeeping.

\section{From Kronecker to Dirac Normalization}

Because \(\Delta k=1/r\), define continuum-normalized states at the sampled
wave numbers by
\begin{equation}
  \ket{k_n}:=\frac{1}{\sqrt{\Delta k}}\,
    \ket{k_n}_{\mathrm{box}}
    =\sqrt r\,\ket{k_n}_{\mathrm{box}}.
  \label{eq:continuum-rescaling}
\end{equation}
Their position-space kernels are
\begin{equation}
  \braket{x}{k_n}
    =\sqrt r\,\frac{e^{ik_nx}}{\sqrt{2\pi r}}
    =\frac{e^{ik_nx}}{\sqrt{2\pi}}.
  \label{eq:continuum-plane-wave}
\end{equation}
The factor of \(\sqrt r\) has not vanished. It has converted a
unit-normalized discrete state into a delta-normalized continuum state.
Indeed,
\begin{equation}
  \braket{k_m}{k_n}
    =\frac{\delta_{mn}}{\Delta k},
  \label{eq:sampled-delta}
\end{equation}
which tends distributionally to
\begin{equation}
  \braket{k'}{k}=\delta(k-k').
  \label{eq:momentum-delta}
\end{equation}
The peak grows like \(1/\Delta k=r\) while its width shrinks like
\(\Delta k\), leaving unit area.

\begin{figure}[tbp]
  \centering
  \includegraphics[width=0.66\textwidth]{lecture_05_figure_04.png}
  \caption{The board transition from normalized plane waves to
  \(\delta(k-k')\) normalization.}
  \lectureframe{00:16:50}
\end{figure}

\begin{classroomqa}
\audiencequestion{Where did the factor \(\sqrt r\) go in the overlap?
\lecturetimestamp{00:17:37--00:18:39}}
\lecturerresponse{There is one factor from the bra and one from the ket.
Together they multiply the old overlap by \(r\). Thus unequal sampled wave
numbers remain orthogonal, while the equal-wave-number overlap grows from
\(1\) to \(r=1/\Delta k\), exactly the height needed for a Dirac delta in
the continuum limit.}
\end{classroomqa}

Completeness transforms consistently:
\begin{align}
  \mathbf 1
    &=\sum_n \ket{k_n}_{\mathrm{box}}
       {}_{\mathrm{box}}\!\bra{k_n}\\
    &=\sum_n \Delta k\,\ket{k_n}\bra{k_n}
      \longrightarrow
      \int_{-\infty}^{\infty}dk\,\ket k\bra k.
  \label{eq:continuum-completeness}
\end{align}
The rescaling, the delta function, and the replacement of a sum by an
integral are three parts of one operation.

\section{Dyads and the Resolution of Identity}

An inner product such as \(\braket ab\) is a number. Reversing the order
creates a different kind of object:
\begin{equation}
  \ket b\bra a.
\end{equation}
This outer product, or dyad, is an operator. Acting on a third vector,
\begin{equation}
  \bigl(\ket b\bra a\bigr)\ket c
    =\braket ac\,\ket b.
  \label{eq:dyad-action}
\end{equation}
The bra extracts the component of \(\ket c\) along \(\ket a\), and the ket
places that coefficient along \(\ket b\).

For any orthonormal basis \(\{\ket n\}\),
\begin{equation}
  \ket v=\sum_n v_n\ket n,
  \qquad
  v_n=\braket n v.
\end{equation}
Substitution gives
\begin{equation}
  \ket v=\sum_n\ket n\bra n\ket v.
\end{equation}
Since the relation holds for every \(\ket v\),
\begin{equation}
  \mathbf 1=\sum_n\ket n\bra n.
  \label{eq:l05-discrete-completeness}
\end{equation}

\begin{figure}[tbp]
  \centering
  \includegraphics[width=0.50\textwidth]{lecture_05_figure_05.png}
  \caption{A complete orthonormal basis resolves the identity into
  projectors.}
  \lectureframe{00:25:02}
\end{figure}

\begin{classroomqa}
\audiencequestion{Does the repeated symbol \(n\) refer to one special
state?\lecturetimestamp{00:25:58--00:27:40}}
\lecturerresponse{No. It is a dummy summation label that runs over every
vector in the chosen basis. Another complete orthonormal basis gives another
resolution of the same identity. For a continuous label, the sum becomes an
integral and Kronecker normalization becomes Dirac normalization.}
\end{classroomqa}

In particular,
\begin{equation}
  \mathbf 1=\int dx\,\ket x\bra x
  =\int dk\,\ket k\bra k.
  \label{eq:two-resolutions}
\end{equation}
These two forms of the identity are the mechanism behind the Fourier
transform.

\section{One State in Two Representations}

For an abstract state \(\ket\psi\), define
\begin{equation}
  \psi(x):=\braket x\psi,
  \qquad
  \widetilde\psi(k):=\braket k\psi.
  \label{eq:two-wavefunctions}
\end{equation}
They are not two states. They are the components of one state in two
different bases. The change-of-basis kernel is
\begin{equation}
  \braket xk=\frac{e^{ikx}}{\sqrt{2\pi}},
  \qquad
  \braket kx=\frac{e^{-ikx}}{\sqrt{2\pi}}.
  \label{eq:fourier-kernel}
\end{equation}

Insert the position resolution of identity into \(\braket k\psi\):
\begin{align}
  \widetilde\psi(k)
    &=\int dx\,\braket kx\braket x\psi\\
    &=\frac{1}{\sqrt{2\pi}}
      \int_{-\infty}^{\infty}dx\,e^{-ikx}\psi(x).
  \label{eq:forward-fourier}
\end{align}
Insert the momentum resolution into \(\braket x\psi\):
\begin{align}
  \psi(x)
    &=\int dk\,\braket xk\braket k\psi\\
    &=\frac{1}{\sqrt{2\pi}}
      \int_{-\infty}^{\infty}dk\,e^{ikx}\widetilde\psi(k).
  \label{eq:inverse-fourier}
\end{align}
These are the forward and inverse Fourier transforms. Fourier analysis is
older than quantum mechanics; quantum mechanics interprets the transform as
a change between position and momentum descriptions.

\begin{figure}[tbp]
  \centering
  \includegraphics[width=0.68\textwidth]{lecture_05_figure_06.png}
  \caption{Inserting the position basis turns the momentum amplitude into a
  Fourier integral.}
  \lectureframe{00:34:30}
\end{figure}

\begin{classroomqa}
\audiencequestion{Can \(k\) be negative, and what does the symmetry between
\(x\) and \(k\) mean physically?\lecturetimestamp{00:38:43--00:39:43}}
\lecturerresponse{Negative \(k\) means momentum in the opposite direction.
The deeper symmetry is that position and momentum are conjugate variables.
Their reciprocity is natural in Fourier analysis and later appears in
classical mechanics through canonical coordinates and Poisson brackets.}
\end{classroomqa}

The operator rules display the same reciprocity. In the position
representation,
\begin{align}
  \hat X\psi(x)&=x\psi(x),\\
  \hat K\psi(x)&=-i\frac{\partial}{\partial x}\psi(x).
  \label{eq:x-representation}
\end{align}
In the wave-number representation,
\begin{align}
  \hat K\widetilde\psi(k)&=k\widetilde\psi(k),\\
  \hat X\widetilde\psi(k)&=
    i\frac{\partial}{\partial k}\widetilde\psi(k).
  \label{eq:k-representation}
\end{align}
The physical momentum operator is \(\hat P=\hbar\hat K\).

\begin{figure}[tbp]
  \centering
  \includegraphics[width=0.62\textwidth]{lecture_05_figure_07.png}
  \caption{Multiplication and differentiation exchange roles in the
  \(x\)- and \(k\)-representations.}
  \lectureframe{00:42:15}
\end{figure}

\begin{classroomqa}
\audiencequestion{Why is the state represented by \(\psi(x)\) or
\(\widetilde\psi(k)\), rather than by one probability amplitude depending
on both \(x\) and \(k\)?\lecturetimestamp{00:43:13--00:44:36}}
\lecturerresponse{These are two complete representations of the same state,
and they are related invertibly by a Fourier transform. Position and
momentum are incompatible observables: no normalizable state has both
exactly sharp. A definite-\(k\) state is spread over all \(x\), while a
definite-position state is spread over all \(k\). Quantum mechanics does not
assign simultaneous sharp values merely because either representation can
be used.}
\end{classroomqa}

\begin{classroomqa}
\audiencequestion{Is the classical conservation of phase-space volume
related to quantum time evolution?\lecturetimestamp{00:44:38--00:46:39}}
\lecturerresponse{Yes, at the structural level emphasized in the lecture.
Hamiltonian evolution is a canonical transformation and preserves
phase-space volume. It cannot collapse a finite cloud of distinct initial
conditions to one point. Quantum evolution is unitary and preserves inner
products, so distinctions between state vectors are likewise maintained.
The detailed dynamical statement is still ahead of us.}
\end{classroomqa}

\begin{classroomqa}
\audiencequestion{Where does the differential form
\(\hat K=-i\partial_x\) come from?\lecturetimestamp{00:46:43--00:48:08}}
\lecturerresponse{It is the position-representation form introduced for the
momentum generator. Acting on a plane wave,
\[
  -i\frac{\partial}{\partial x}e^{ikx}=k e^{ikx},
\]
so \(e^{ikx}\) is an eigenfunction with eigenvalue \(k\). Fourier
transformation then gives the reciprocal rule
\(\hat X=i\partial_k\) in momentum space.}
\end{classroomqa}

\section{Wave Packets and the Classical Bridge}

At this stage, the operators have been named position and momentum, but the
connection with classical motion has not yet been proved. The bridge is a
wave packet. Its magnitude is concentrated near some position \(x_0\), and
its Fourier transform is concentrated near some wave number \(k_0\). A
typical form is
\begin{equation}
  \psi(x)=f(x-x_0)e^{ik_0x},
  \label{eq:wave-packet}
\end{equation}
where \(f\) is a smooth localized envelope. The oscillations inside the
envelope carry the approximate momentum, while the envelope locates the
particle.

\begin{figure}[tbp]
  \centering
  \includegraphics[width=0.68\textwidth]{lecture_05_figure_08.png}
  \caption{A localized envelope modulates a nearly monochromatic plane wave,
  producing a packet localized in both representations.}
  \lectureframe{00:51:25}
\end{figure}

\begin{classroomqa}
\audiencequestion{Must the packet be Gaussian?
\lecturetimestamp{00:50:23--00:50:34}}
\lecturerresponse{No. A Gaussian is the cleanest example because its Fourier
transform is also Gaussian and it can minimize
\(\Delta x\,\Delta k\geq 1/2\). The lecture needs only a smooth packet
localized in both representations; its exact shape is not essential.}
\end{classroomqa}

Under a sufficiently smooth potential, a narrow packet can retain its
integrity for a useful interval. Its center moves through position space,
and the center of its momentum distribution moves through momentum space.
The theorem promised in the lecture is the Ehrenfest relation: under the
conditions in which the packet remains narrow, these expectation values
approximately obey the classical equations of motion. The packet may also
spread or distort, so the classical description is an approximation, not a
new postulate.

\begin{classroomqa}
\audiencequestion{What establishes that this quantum quantity really is
classical momentum?\lecturetimestamp{00:48:13--00:54:02}}
\lecturerresponse{Not the name alone. The connection is established when a
localized wave packet evolves according to the Schrodinger equation and its
centers in \(x\) and \(p=\hbar k\) follow the corresponding classical
trajectory. That result is announced here and derived only after quantum
dynamics has been introduced.}
\end{classroomqa}

\section{Incompatibility in a Two-State System}

For complex vectors, \(\ket A\bra B\) is generally not Hermitian, and
\begin{equation}
  \bigl(\ket A\bra B\bigr)^\dagger=\ket B\bra A.
  \label{eq:dyad-adjoint}
\end{equation}
Hermitian conjugation is the operator analogue of complex conjugation.

\begin{classroomqa}
\audiencequestion{What is the outer product of two complex vectors?
\lecturetimestamp{00:54:15--00:55:32}}
\lecturerresponse{It is an operator, usually a non-Hermitian one. Reversing
the vectors and converting each ket to its adjoint bra gives its Hermitian
conjugate. The special dyad \(\ket A\bra A\) is Hermitian and, for normalized
\(\ket A\), is a projector.}
\end{classroomqa}

The larger point is incompatibility. Classically one specifies position
\textit{and} momentum. Quantum mechanically one may use the position
representation \textit{or} the momentum representation. The state contains
the probability information needed for either experiment, but it does not
make incompatible quantities simultaneously definite.

To see the same logic without infinite-dimensional function spaces, turn to
photon polarization. Ignore the photon's position and retain only two
polarization degrees of freedom.

\section{From Electromagnetic Waves to a Polarizer}

A classical light wave travelling along the \(z\)-axis has electric and
magnetic fields transverse to its motion and perpendicular to one another.
For a plane wave they oscillate in phase. The polarization direction is
defined to be the direction in which the electric field oscillates; the
opposite direction along the same axis is not a different linear
polarization.

\begin{figure}[tbp]
  \centering
  \includegraphics[width=0.68\textwidth]{lecture_05_figure_09.png}
  \caption{Perpendicular electric and magnetic oscillations in a
  plane-polarized electromagnetic wave.}
  \lectureframe{01:00:45}
\end{figure}

A wire-grid polarizer makes this physical. If its conducting wires are
vertical, a vertical electric field drives currents in the wires. Those
currents reflect or absorb that component. A horizontal electric field
cannot drive a horizontal current because no conducting path runs that way,
so it is transmitted. Thus the transmission axis is perpendicular to the
wires. More generally, an ideal polarizer transmits the electric-field
component along one axis and removes the orthogonal component.

The wave description gives continuously variable field amplitudes. The
single-photon description gives discrete events. A photon incident on a
polarizer is either transmitted or it is absorbed or reflected. Its energy
does not split into two fractional photons.

\section{Preparation and Measurement}

A photon that emerges from an ideal vertical polarizer is vertically
polarized. If it then meets a second vertical polarizer, it passes with
probability one. If it meets a horizontal polarizer, it is blocked with
probability one. The first apparatus prepares a state; the second tests it.
The same physical device can therefore serve as a state preparer or as a
measurement apparatus, depending on its role in the experiment.

\begin{figure}[tbp]
  \centering
  \includegraphics[width=0.66\textwidth]{lecture_05_figure_10.png}
  \caption{Two polarizers illustrate preparation followed by a polarization
  measurement.}
  \lectureframe{01:10:00}
\end{figure}

\begin{classroomqa}
\audiencequestion{Is a reflected photon itself polarized?
\lecturetimestamp{01:06:14--01:06:38}}
\lecturerresponse{It may be. The lecture discards that output channel and
counts only transmission. One may instead build a fuller experiment that
detects both channels, but the two-outcome polarization logic for the
transmitted beam is unchanged.}
\end{classroomqa}

Choose horizontal and vertical axes called \(x\) and \(y\). In this section
they label polarization, not the particle's position. Represent the two
orthogonal states by
\begin{equation}
  \ket{x}_{\mathrm{pol}}\leftrightarrow
  \begin{pmatrix}1\\0\end{pmatrix},
  \qquad
  \ket{y}_{\mathrm{pol}}\leftrightarrow
  \begin{pmatrix}0\\1\end{pmatrix}.
  \label{eq:xy-polarization-basis}
\end{equation}
They satisfy
\begin{equation}
  \braket xx=\braket yy=1,
  \qquad
  \braket xy=0.
  \label{eq:xy-orthogonality}
\end{equation}
They are orthogonal because a single ideal polarization experiment can
distinguish them perfectly.

\section{States, Outcomes, and the Polarization Operator}

The state vectors are not themselves observables. To define an observable,
associate numerical records with the two possible outcomes. Following the
lecture, record \(+1\) for \(x\)-polarization and \(-1\) for
\(y\)-polarization. The corresponding operator obeys
\begin{equation}
  \hat P_{xy}\ket x=+\ket x,
  \qquad
  \hat P_{xy}\ket y=-\ket y.
  \label{eq:Pxy-eigenvalue}
\end{equation}
In the \((\ket x,\ket y)\) basis,
\begin{equation}
  \hat P_{xy}
  =
  \begin{pmatrix}
    1&0\\
    0&-1
  \end{pmatrix}
  =\ket x\bra x-\ket y\bra y.
  \label{eq:Pxy-matrix}
\end{equation}

\begin{figure}[tbp]
  \centering
  \includegraphics[width=0.58\textwidth]{lecture_05_figure_11.png}
  \caption{The horizontal-vertical polarization observable is diagonal in
  the \((\ket x,\ket y)\) basis.}
  \lectureframe{01:19:45}
\end{figure}

\begin{classroomqa}
\audiencequestion{If polarization is observable, why is the state not called
an observable?\lecturetimestamp{01:14:56--01:16:03}}
\lecturerresponse{The state describes the preparation and supplies
probability amplitudes. The observable is the Hermitian operator tied to a
measurement setup; its eigenvectors are states with definite outcomes, and
its eigenvalues are the numbers recorded for those outcomes.}
\end{classroomqa}

Knowing the action of \(\hat P_{xy}\) on the two basis vectors determines its
action on every superposition. This is why specifying eigenvectors and
eigenvalues is enough to specify an observable.

\section{Rotate the Apparatus}

Nothing prevents us from rotating a polarizer by \(45^\circ\). A photon that
emerges from it is neither in \(\ket x\) nor in \(\ket y\), but an
\(x\)- or \(y\)-oriented test gives either result with equal probability.
The normalized state is
\begin{equation}
  \ket{+}=\frac{\ket x+\ket y}{\sqrt2}.
  \label{eq:plus-state}
\end{equation}
The orthogonal linear polarization is
\begin{equation}
  \ket{-}=\frac{\ket x-\ket y}{\sqrt2}.
  \label{eq:minus-state}
\end{equation}

\begin{figure}[tbp]
  \centering
  \includegraphics[width=0.66\textwidth]{lecture_05_figure_12.png}
  \caption{A diagonal preparation followed by a horizontal or vertical
  polarization test.}
  \lectureframe{01:23:45}
\end{figure}

The relative minus sign makes the two diagonal states orthogonal:
\begin{equation}
  \braket{-}{+}
  =\frac12\bigl(\braket xx+\braket xy-\braket yx-\braket yy\bigr)
  =0.
  \label{eq:diagonal-orthogonality}
\end{equation}
In column form,
\begin{equation}
  \ket+\leftrightarrow
  \frac{1}{\sqrt2}\begin{pmatrix}1\\1\end{pmatrix},
  \qquad
  \ket-\leftrightarrow
  \frac{1}{\sqrt2}\begin{pmatrix}1\\-1\end{pmatrix}.
  \label{eq:diagonal-columns}
\end{equation}

\begin{figure}[tbp]
  \centering
  \includegraphics[width=0.66\textwidth]{lecture_05_figure_13.png}
  \caption{The two diagonal polarization states form a second orthonormal
  basis.}
  \lectureframe{01:29:00}
\end{figure}

The Born rule confirms the equal probabilities:
\begin{align}
  \operatorname{Prob}(y|+)
    &=\left|\braket y+\right|^2
      =\left|\frac{1}{\sqrt2}
        \bigl(\braket yx+\braket yy\bigr)\right|^2\\
    &=\frac12,\\
  \operatorname{Prob}(x|+)&=\frac12.
  \label{eq:diagonal-probabilities}
\end{align}
For \(\ket-\), one amplitude changes sign, but its absolute square is again
\(1/2\). Amplitudes carry signs and phases; probabilities do not retain
them directly.

Define the rotated observable by
\begin{equation}
  \hat P_{45}\ket+=+\ket+,
  \qquad
  \hat P_{45}\ket-=-\ket-.
  \label{eq:P45-eigenvalue}
\end{equation}
In the original \(x,y\) basis,
\begin{equation}
  \hat P_{45}
  =
  \begin{pmatrix}
    0&1\\
    1&0
  \end{pmatrix}
  =\ket+\bra+-\ket-\bra-.
  \label{eq:P45-matrix}
\end{equation}

\begin{figure}[tbp]
  \centering
  \includegraphics[width=0.62\textwidth]{lecture_05_figure_14.png}
  \caption{The off-diagonal matrix returns \(\ket+\) and reverses
  \(\ket-\), giving the required eigenvalues.}
  \lectureframe{01:38:30}
\end{figure}

The eigenvectors of \(\hat P_{xy}\) are \(\ket x,\ket y\); those of
\(\hat P_{45}\) are \(\ket+,\ket-\). No nonzero vector is an eigenvector of
both operators. Equivalently,
\begin{equation}
  [\hat P_{xy},\hat P_{45}]\neq0.
  \label{eq:polarization-commutator}
\end{equation}
The two polarization observables are incompatible, just as position and
momentum are incompatible, but now the entire statement fits inside a
two-dimensional vector space.

For an arbitrary linear-polarization angle \(\theta\), the announced
generalization is
\begin{align}
  \ket{\theta}
  &=\cos\theta\,\ket x+\sin\theta\,\ket y,\\
  \ket{\theta_\perp}
  &=-\sin\theta\,\ket x+\cos\theta\,\ket y.
  \label{eq:general-linear-polarization}
\end{align}
The corresponding observable has these two states as eigenvectors. Rotating
the axis by \(90^\circ\) merely exchanges the outcomes and reverses the sign
of the same observable. Distinct nonorthogonal axes give genuinely
incompatible tests. Complex relative coefficients introduce circular and
elliptical polarization; that extension is reserved for the next lecture.

\section{What One Experiment Records}

The closing discussion distinguishes an outcome from a probability. A
single photon gives one definite record: transmitted or blocked, represented
here by \(+1\) or \(-1\). The probability of either record is inferred from
many repetitions with identically prepared photons.

\begin{classroomqa}
\audiencequestion{Is the probability that a photon passes a polarizer itself
an observable?\lecturetimestamp{01:42:27--01:45:03}}
\lecturerresponse{Not in the single-shot technical sense used here. One run
produces an outcome, not a probability. Probability is estimated from the
relative frequency of outcomes over an ensemble of repeated preparations.}
\end{classroomqa}

The probability for a particular eigenstate \(\ket a\) is obtained directly
from the overlap,
\begin{equation}
  \operatorname{Prob}(a|\psi)=\left|\braket a\psi\right|^2.
  \label{eq:born-overlap}
\end{equation}
One does not first apply the observable to \(\ket\psi\) and then square an
eigenvalue. The observable identifies the possible values and the
eigenvectors associated with them; the state's overlaps with those
eigenvectors determine the probabilities.

\begin{classroomqa}
\audiencequestion{Should the observable act on the state before the
probability is calculated?\lecturetimestamp{01:45:03--01:48:02}}
\lecturerresponse{No. To find the probability of a particular result, take
the inner product of the prepared state with the eigenvector for that
result, then square its magnitude. Applying the operator is useful for
expectation values and eigenvalue equations, but it is not an extra step in
the basic Born-rule calculation.}
\end{classroomqa}

\begin{classroomqa}
\audiencequestion{Is an observable synonymous with a Hermitian operator?
\lecturetimestamp{01:48:04--01:48:14}}
\lecturerresponse{Within the finite-dimensional framework of this lecture,
yes: observables are represented by Hermitian operators, and every
Hermitian operator may be associated with an observable. In
infinite-dimensional problems, the precise mathematical condition is
self-adjointness with a specified domain.}
\end{classroomqa}

The numerical labels \(+1\) and \(-1\) are a convention. We could record
\(1\) and \(0\), or any two distinct real numbers, provided the convention
is used consistently. Choosing \(1\) and \(0\) produces the projector
\begin{equation}
  \Pi_x=\ket x\bra x
  =
  \begin{pmatrix}
    1&0\\
    0&0
  \end{pmatrix}.
  \label{eq:x-projector}
\end{equation}
For a normalized state \(\ket\psi\),
\begin{align}
  \bra\psi\Pi_x\ket\psi
    &=\bra\psi x\rangle\braket x\psi\\
    &=\left|\braket x\psi\right|^2.
  \label{eq:projector-probability}
\end{align}
The expectation value equals the probability, but the result of one
projector measurement is still \(1\) or \(0\).

\begin{figure}[tbp]
  \centering
  \includegraphics[width=0.42\textwidth]{lecture_05_figure_15.png}
  \caption{A projector expectation gives the \(x\)-polarization
  probability.}
  \lectureframe{01:50:30}
\end{figure}

\begin{classroomqa}
\audiencequestion{If the eigenvalue labels are arbitrary, what does the
apparatus actually define?\lecturetimestamp{01:48:16--01:55:23}}
\lecturerresponse{The apparatus defines the mutually exclusive outcomes and
the states that give them with certainty. We then choose real numbers to
record those outcomes. With the \(+1,-1\) convention the observable is a
difference of projectors; with the \(1,0\) convention it is one projector,
whose expectation value is the ensemble probability for that outcome.}
\end{classroomqa}

An observable is therefore inseparable from a specified experimental
question. The polarizer orientation fixes the eigenvectors, the notebook
convention fixes the eigenvalues, and one interaction yields one recorded
event. Repetition reveals the probabilities. In this two-state example, the
abstract language of vectors, operators, incompatible bases, and Born
amplitudes has become a concrete laboratory procedure.
